\section{The five properties and their stories}
\label{sec:five-properties}
The following meanings state the adopted design in plain words. Checked
results concern the specified model; application behavior and a running
demonstration require evidence at their own boundary. Every checked reference
uses the exact revision and axiom line in Appendix~\ref{sec:artifact}.

\subsection{Named}
Each token has a name: registering a name mints exactly one token for it,
and the registry is the only way to mint it. Holding the token is holding
the name. The concrete story is that a registered name has a token which
an application can hold and associate with its own payload.

\checked The original registry's accepted registration transaction mints
one active asset indexed by the request's key, delivers it to the named
destination and records that holding. Its model asset name equals the key.
\evidence{insert_active_transaction_row}
\checked In an economic reachable world, actual singular-tag supply across
application custody and pending terminations is one exactly for an active
key, and zero otherwise.
\evidence{world_supply_sync}
\missing These results are about modeled identities and transitions.
``Only way to mint'' on Cardano requires correspondence of the actual minting
policy and name encoding. Holding the name's representative is not a proof
of a legal right, personal identity or entitlement to that name. The
application decides what its payload means.

\subsection{Unique}
A name can be taken only with a proof that it is absent, so two people cannot
hold separate genuine tokens for the same name in one registry. Bob tries
to register Alice's name after she has registered it; his registration is
refused.

\checked After an accepted original registration, another registration of
the same key with valid approval is refused by the modeled builder with
\texttt{key-exists}. The same declaration records exactly one active holding.
\evidence{insert_active_transaction_row}
\checked The economic supply invariant retains the bound of one while the
token moves from application custody into a pending termination.
\evidence{world_supply_sync}
\missing A concrete absence proof's verification by the validator is a
correspondence obligation. ``Absent'' here means an unknown, untaken key,
not the historical stored Absent leaf. The refusal concerns taking an
already active name; booking a competing request does not reserve it, and
the model does not adjudicate who ought to have registered it first.

\subsection{Permanent}
Terminating a name burns its token and records it as terminal; it can never
be issued again. Alice terminates her name and then tries to register it
again. The attempt is refused. For KERI, termination is intended to be
conviction, so a convicted identifier must never come back in that registry.

\checked For a terminal key in a reachable original registry, every request
at that key is refused; every accepted later batch keeps it terminal.
\evidence{termination}
\checked A genuine pending termination has an accepted economic fold that
makes its key terminal under the theorem's explicit funding and willingness
premises. World supply is then zero for that key.
\evidence{booked_termination_foldable}
\evidence{world_supply_sync}
\missing The application decides when release from custody is authorized.
A KERI conviction needs its own duplicity-proof and integration evidence.
Permanence within this registry does not identify a new registry with the old
one. Economically, the termination donates everything held except the tip
to the treasury; it returns no deposit to Alice.
\evidence{fold_routes_every_amount}

\subsection{Permissionless}
Anyone creates a registry, anyone books a name, anyone settles pending
requests, all from public chain data. The creator has no power over the
names. These are the design's participation promises, subject to application
authorization, funding and the admission rules. The concrete story is Bob,
who did not create the registry, settling the pending requests, Alice's included.

\checked The original registry builder is equal to a construction from its
public view: the registry and pending request, including the request's datum.
An accepted original fold requires no signer, and changing approval signatures
does not change that step or its built transaction.
\evidence{fold_inputs_public}
\evidence{fold_requires_no_signer}
\missing These are original-builder statements, not a public-data theorem
for the economic transaction builder or an execution of Bob's story. Registry
creation, obtaining all chain data, funding and submitting the actual batch
require off-chain and validator evidence. The story is a demonstration of
that boundary, not a model theorem or an inclusion guarantee. Creator
privilege over application payloads is likewise not decided by these results.

\subsection{Reusable}
The same mint serves any application: open-datum today, a closed variant in
the demonstration, KERI checkpoints next, with no change to its contract.
Each name's token lives in its owner's application, which decides what the
name points to. This is the design's separation of the name from its payload.

\missing The concrete application demonstration uses a closed flag on the
same release and commands: a payload set at registration cannot be updated,
the application refuses every update, and only termination remains. This
is application-validator and executable behavior outside the checked registry
model. No release or running-demonstration receipt is asserted by this paper.
KERI is the next intended consumer, not a completed integration. Serving
both application kinds within one registry through application tags is also
a future integration claim. The application's release rule and adaptations
to the ruled termination economics remain open.
